\tikzset{
  % --- ESTILOS PARA NODOS ---
  % Nodo base: un punto negro. El tamaño se controla con 'inner sep'.
  nodo/.style = {
    circle,
    fill=black,
    inner sep=1.5pt
  },
  % Nodo resaltado: es un nodo normal con un borde grueso y gris que actúa como halo.
  nodo_resaltado/.style = {
    nodo,
    draw=gray!40,
    line width=4pt
  },
  % Estilo para las etiquetas de los nodos (que van por fuera).
  etiqueta_nodo/.style = {
    font=\small,
    text=black
  },
  % Alias de compatibilidad para figuras antiguas.
  etiqueta/.style = {
    etiqueta_nodo
  },
  % --- ESTILOS PARA ARISTAS ---
  % Arista base: línea negra y gruesa.
  arista/.style = {
    thick,
    black
  },
  % Flecha: hereda de 'arista' y añade una punta.
  flecha/.style = {
    arista,
    ->,
    >=Stealth
  },
  % Arista resaltada: usa 'preaction' para dibujar una línea gris más ancha por debajo.
  arista_resaltada/.style = {
    arista,
    preaction={
        draw=gray!40,
        line width=4pt
    }
  },
  % Flecha resaltada: igual que la arista, pero el halo también es una flecha.
  flecha_resaltada/.style = {
    flecha,
    preaction={
        draw=gray!40,
        line width=4pt,
        ->, >=Stealth
    }
  },
  % Arista/Flecha tenue: para elementos de fondo.
  arista_tenue/.style = {
    arista,
    draw=gray!50
  },
  flecha_tenue/.style = {
    flecha,
    draw=gray!50
  },
  % Flecha alternativa: usada para rutas opcionales en ejemplos.
  flecha_alternativa/.style = {
    flecha_tenue,
    dashed
  },
  % --- ESTILO PARA ETIQUETAS DE PESO ---
  etiqueta_peso/.style = {
    midway,
    sloped,
    fill=white,
    inner sep=1.5pt,
    font=\small
  }
}
